% ===== PRÉAMBULE CANONIQUE — à coller VERBATIM en tête de CHAQUE .tex =====
\documentclass[11pt,a4paper]{article}
\usepackage[utf8]{inputenc}\usepackage[T1]{fontenc}\usepackage{lmodern}
\usepackage[french]{babel}
\usepackage{amsmath,amssymb,mathtools}\usepackage{icomma}
\usepackage{siunitx}\sisetup{locale=FR}  % \qty{}{}, \si{}, \ang{}
\usepackage{xcolor}\usepackage{colortbl}
\usepackage{tikz}\usepackage{pgfplots}\pgfplotsset{compat=1.18}
\usepackage[siunitx,europeanresistors]{circuitikz} % circuits (élec.)
\usepackage[most]{tcolorbox}\usepackage{enumitem}\usepackage{array,booktabs}
\usepackage[a4paper,margin=2.2cm]{geometry}
\usetikzlibrary{arrows.meta,calc,angles,quotes,positioning,patterns,%
  backgrounds,shapes.geometric,decorations.pathmorphing,decorations.markings,intersections}
\definecolor{primary}{RGB}{222,49,99}\definecolor{primarydark}{RGB}{150,22,55}
\definecolor{defcol}{RGB}{0,114,178}\definecolor{methcol}{RGB}{52,92,148}
\definecolor{excol}{RGB}{0,135,90}\definecolor{attncol}{RGB}{200,40,55}
\tcbset{boxbase/.style={enhanced,breakable,boxrule=0.9pt,arc=2.5pt,
  left=3mm,right=3mm,top=2.3mm,bottom=2.3mm,fonttitle=\bfseries,coltitle=white}}
% --- boîtes TUTORIEL (noms EXACTS, ne pas renommer) ---
\newtcolorbox{cle}[1][]{boxbase,colback=defcol!6,colframe=defcol,
  title={\textsc{À retenir}\ifx\relax#1\relax\relax\else\textmd{ : \itshape #1}\fi}}
\newtcolorbox{methode}[1][]{boxbase,colback=methcol!7,colframe=methcol,sharp corners,
  title={\textsc{Méthode}\ifx\relax#1\relax\relax\else\textmd{ --- \itshape #1}\fi}}
\newtcolorbox{exemplebox}[1][]{boxbase,colback=excol!5,colframe=excol,
  title={\textsc{Exemple}\ifx\relax#1\relax\relax\else\textmd{ : \itshape #1}\fi}}
\newtcolorbox{piege}[1][]{boxbase,colback=attncol!6,colframe=attncol,
  title={\textsc{Piège}\ifx\relax#1\relax\relax\else\textmd{ : \itshape #1}\fi}}
\newtcolorbox{remarque}[1][]{boxbase,colback=black!4,colframe=black!45,
  title={\textsc{Remarque}\ifx\relax#1\relax\relax\else\textmd{ : \itshape #1}\fi}}
% --- boîtes EXERCICE / CORRIGÉ (noms EXACTS, auto-comptées) ---
\newtcolorbox[auto counter]{exo}[1][]{boxbase,colback=primary!4,colframe=primary,
  title={\textsc{Exercice}~\thetcbcounter\ifx\relax#1\relax\relax\else\textmd{ --- \itshape #1}\fi}}
\newtcolorbox[auto counter]{corrige}[1][]{boxbase,colback=excol!4,colframe=excol,
  title={\textsc{Corrigé}~\thetcbcounter\ifx\relax#1\relax\relax\else\textmd{ --- \itshape #1}\fi}}
\newcommand{\vv}[1]{\vec{#1}}\newcommand{\ds}{\displaystyle}
\newcommand{\R}{\mathbb{R}}\newcommand{\N}{\mathbb{N}}\newcommand{\Z}{\mathbb{Z}}
% (un \usepackage{fancyhdr}+en-tête est possible ; il est ignoré côté web)
% =========================================================================
\begin{document}
\sisetup{per-mode=symbol}

\begin{corrige}[conversions d'unités]
\begin{enumerate}[label=\alph*)]
  \item $\ds \qty{72}{\kilo\metre\per\hour}=\frac{72}{3{,}6}=\qty{20}{\metre\per\second}$.
  \item $\ds \qty{30}{\metre\per\second}=30\times3{,}6=\qty{108}{\kilo\metre\per\hour}$.
\end{enumerate}
\end{corrige}

\begin{corrige}[vitesse moyenne et signe]
\begin{enumerate}[label=\alph*)]
  \item $\ds v_{\text{moy}}=\frac{x_2-x_1}{\Delta t}=\frac{170-20}{10}=\qty{15}{\metre\per\second}$ (positive : déplacement dans le sens de l'axe).
  \item $\ds v_{\text{moy}}=\frac{70-170}{5}=\frac{-100}{5}=\qty{-20}{\metre\per\second}$. Le signe négatif indique un déplacement \emph{à contre-sens} de l'axe (le mobile revient).
\end{enumerate}
\end{corrige}

\begin{corrige}[accélération moyenne]
\begin{enumerate}[label=\alph*)]
  \item $\ds a_{\text{moy}}=\frac{25-10}{5}=\qty{3}{\metre\per\second\squared}$.
  \item $\ds a_{\text{moy}}=\frac{0-30}{6}=\qty{-5}{\metre\per\second\squared}$ : accélération négative (freinage).
\end{enumerate}
\end{corrige}

\begin{corrige}[distance parcourue et déplacement]
\begin{enumerate}[label=\alph*)]
  \item Position initiale $0$, finale \qty{10}{\metre} : $\Delta x=10-0=\qty{10}{\metre}$.
  \item Aller $0\to\qty{40}{\metre}$ = \qty{40}{\metre} ; retour $40\to\qty{10}{\metre}$ = \qty{30}{\metre}. Distance totale $=40+30=\qty{70}{\metre}$.
  \item Durée totale $8+6=\qty{14}{\second}$ : $\ds v_{\text{moy}}=\frac{\Delta x}{\Delta t}=\frac{10}{14}\approx\qty{0{,}71}{\metre\per\second}$. (Attention : on divise le \emph{déplacement}, pas la distance.)
\end{enumerate}
\end{corrige}

\begin{corrige}[dériver une équation horaire]
On dérive $x(t)=2t^2+5t+1$.
\begin{enumerate}[label=\alph*)]
  \item $v(t)=\dfrac{\mathrm{d}x}{\mathrm{d}t}=4t+5$ (en \si{\metre\per\second}).
  \item $a(t)=\dfrac{\mathrm{d}v}{\mathrm{d}t}=4\ \si{\metre\per\second\squared}$ (constante $\Rightarrow$ mouvement uniformément accéléré).
  \item $v(3)=4\times3+5=\qty{17}{\metre\per\second}$.
\end{enumerate}
\end{corrige}

\end{document}
